\section{Wprowadzenie i oznaczenia matematyczne}

\begin{frame}{Motywacja}
   Moje zainteresowania naukowe koncentrują się wokół uczenia maszynowego ze wzmocnieniem (RL), dziedziny, w której fundamentalną rolę odgrywają Markowskie Procesy Decyzyjne (MDP) oraz algorytmy wyznaczania polityk optymalnych. Niniejsza praca poświęcona jest analizie modeli MDP z dyskontowaniem quasi-hiperbolicznym, które wprowadza niestacjonarność preferencji w czasie. Takie podejście pozwala lepiej modelować ludzkie zachowania, w szczególności silną preferencję do otrzymywania nagród natychmiastowych. W przyszłości może to posłużyć do tworzenia bardziej ,,ludzkich'' agentów AI m.in. w grach losowych.
\end{frame}

\begin{frame}{Oznaczenia matematyczne}
  \begin{itemize}
    \item $S$ - zbiór stanów
  \item $A$ - zbiór akcji
  \item $q$ - prawdopodobieństwo przejścia $q(s'|s,a)$
  \item $r$ - $S \times A \rightarrow R$, funkcja nagrody
  \item $\alpha > 0$ - parametr skłonności do teraźniejszości
  \item $\beta \in [0,1)$ - standardowy współczynnik dyskontowania wykładniczego
  \item $\Phi$, $\Phi_m$, $\Phi_s$ -- odpowiednio polityki zależne od historii, markowskie i stacjonarne.
  \end{itemize}
\end{frame}

